\icode{E-MCC-MU-L3-261025-001}
（多选）
如图所示的皮带传动装置，主动轮\(O_1\)上两轮的半径分别为\(3r\)和\(r\)，从动轮\(O_2\)的半径为\(2r\)，\(A\)、\(B\)、\(C\)分别为轮缘上的三点，设皮带不打滑，则下比例正确的是\paren[]

\begin{tikzpicture}[font=\scriptsize]
	\draw (0,0) circle(1.5) (0,0) circle(.5) (3,0) circle(1);
	\draw (99.6:.5)--($(3,0)+(99.6:1)$) (-99.6:.5)--($(3,0)+(-99.6:1)$);
	\foreach \c/\str/\tw in {(0,0)/O_1/below,(3,0)/O_2/below,(150:1.5)/A/left,(.5,0)/B/right,(4,0)/C/right}
		\fill \c node[\tw]{\(\str\)} circle(.04);
\end{tikzpicture}
\begin{choices}
	\item \(A\)、\(B\)、\(C\)三点的加速度之比\( a_A:a_B:a_C =6:2:1\)
	\item \(A\)、\(B\)、\(C\)三点的线速度之比\(v_A:v_B:v_C = 3:1:1\)
	\item \(A\)、\(B\)、\(C\)三点的角速度之比\(\omega_A:\omega_B:\omega_C = 2:2:1\)
	\item \(A\)、\(B\)、\(C\)三点的加速度之比\(a_A:a_B:a_C = 3:2:1\)
\end{choices}
\begin{answer}
	ABC\qquad 提示：
\end{answer}